\begin{afterword}
\summaryhead

The post opens with a parable. As Socrates drinks the hemlock, a student of philosophy assures an onlooker that Socrates must die, since all men are mortal ``by definition.'' When Socrates does not die, the student concludes, with the same ``logical certainty,'' that Socrates is not a man.

The author then argues a dilemma. If ``all humans are mortal'' is a generalization from observed cases, the prediction that Socrates will die is a good guess but not certain. If mortality is part of the definition of ``human,'' the conclusion is certain, but you cannot know that Socrates is human until you have seen him die.

More generally, a valid syllogism is valid in every possible world, so observing its validity is no evidence about which world we live in. Logic still has a use: it shows what our uncertain premises predict, and can settle what we did not yet know follows from what we accept. It never settles an empirical question by itself.

\responsehead

The post's argument is a dilemma, and both horns hold. Either the major premise is an empirical generalization, and the conclusion inherits its uncertainty, or it is a definition, and the minor premise, ``Socrates is a man,'' can no longer be known without the conclusion. The parable shows the second horn well: the student keeps the certainty by giving up any prediction.

The dilemma is old, and the post does not say whose it is. John Stuart Mill set it out in \textsc{A System of Logic} (1843), with the same syllogism. He granted that every syllogism, taken as proof, contains a petitio principii. He noted that if ``man'' included mortality, the minor premise would already assert the conclusion. And his answer, that ``All inference is from particulars to particulars,'' is the post's own alternative: the last thirty humans who drank hemlock died, so Socrates will. The earliest printing of the example itself that one researcher, David Wheeler, could find is in Mill's chapter, which fits the question mark the post puts after ``Aristotle.'' Mill's authority supports the post; a reader should know that the argument has it.

The historical foil is drawn loosely. The claim that Greek philosophers put mortality into the definition of man has a basis: Porphyry's definition of man as a rational mortal animal was adopted by medieval logicians. The claim that they were ``rather fond of certainty'' is given without support and does not fit every school; Plato's own Academy turned skeptical under Arcesilaus and Carneades.

The second half is correct. Evidence favors a hypothesis only if it is likelier when the hypothesis is true, and the validity of a form holds in every world, so it favors nothing. This is the textbook difference between a valid argument and a sound one, stated in terms of evidence. The wording ``syllogisms \ldots\ are always valid'' is loose, since most syllogistic forms are invalid, but the sense is clear. The post also limits its own claim: logic can tell you whether 29384209 is prime. The idea it borrows for this, ``impossible possible worlds,'' is credited only as ``It's been suggested''; the phrase comes from the title of a 1975 paper by Jaakko Hintikka.

In short: a sound dilemma about truth ``by definition'' and a correct account of what validity cannot show; the dilemma, with the same syllogism and the same inductive alternative, is John Stuart Mill's, from 1843.
\end{afterword}
